% Hauptdatei des gesamten Messprotokolls.
% Immer diese Datei kompilieren; die Kapitel werden unten eingebunden.

\documentclass{include/berichtclass}
\usepackage{wrapfig}
\usepackage[numbers]{natbib}
% !TeX root = ../main.tex
% Eigene Hilfsbefehle und mathematische Kurzschreibweisen.
% Diese Datei wird bereits in main.tex eingebunden.

\newcommand{\op}[1]{\operatorname{#1}}
\newcommand{\arr}[2]{\begin{array}{#1}#2\end{array}}
\renewcommand{\d}{\ensuremath{\text{d}}}

\newcommand{\NN}{\mathbb{N}}
\newcommand{\RR}{\mathbb{R}}
\newcommand{\CC}{\mathbb{C}}
\newcommand{\pdb}[2]{\frac{\partial #1}{\partial #2}}

% Markierung offener Stellen im Protokoll: \todo{Messwert prüfen}
\newcommand{\margtodo}{\marginpar{\textbf{\textcolor{expcolor}{ToDo}}}}
\newcommand{\todo}[1]{{\textbf{\textcolor{expcolor}{[\margtodo{}#1]}}}}
\newenvironment{deprecated}{\begin{color}{gray}}{\end{color}}

% Abbildungs-/Tabellenhilfen aus der ursprünglichen Vorlage
\newcommand{\xcaption}[2]{\caption[#1]{\textbf{#1} #2}}
\newcommand{\xfigure}[5]{%
    \begin{figure}[#1]
        \centering
        \includegraphics[scale=#2]{./fig/#3}
        \xcaption{#4}{#5}
        \label{fig:#3}
    \end{figure}}
\newcommand{\xfigurerot}[5]{%
    \begin{figure}[#1]
        \centering
        \includegraphics[angle=270,scale=#2]{./fig/#3}
        \xcaption{#4}{#5}
        \label{fig:#3}
    \end{figure}}
\newcommand{\xtable}[4]{%
    \begin{table}[#1]
        \centering
        \xcaption{#3}{#4}
        \rowcolors{3}{gray!10}{white}
        \include{./tab/#2}
        \label{tab:#2}
    \end{table}}

% Quantenmechanik und Mathematik
\newcommand{\ket}[1]{\left|#1\right\rangle}
\newcommand{\bra}[1]{\left\langle#1\right|}
\newcommand{\braket}[2]{\left\langle#1 \middle| #2\right\rangle}
\newcommand{\bratenket}[3]{\left\langle#1 \middle|\middle| #2 \middle|\middle| #3\right\rangle}
\newcommand{\anglemean}[1]{\left\langle #1 \right\rangle}
\newcommand{\norm}[1]{\left\lVert #1 \right\rVert}

\newcommand{\updownarrows}{\text{\rotatebox[origin=c]{90}{$\rightleftarrows$}}}
\newcommand{\downuparrows}{\text{\rotatebox[origin=c]{270}{$\rightleftarrows$}}}
\newcommand{\neswarrows}{\text{\rotatebox[origin=c]{45}{$\rightleftarrows$}}}
\newcommand{\swnearrows}{\text{\rotatebox[origin=c]{225}{$\rightleftarrows$}}}

\newcommand{\cre}{c^\dagger}
\newcommand{\anh}{c^{\vphantom{\dagger}}}
\newcommand{\numb}{n^{\vphantom{\dagger}}}
\newcommand{\fullstop}{\text{\,.}}
\newcommand{\comma}{\text{\,,}}


\SelectLanguage{ngerman}

% --------------------------------------------------
% Einstellungen für diesen Versuch: hier anpassen
% --------------------------------------------------
\newcommand{\versuch}{Versuch XX -- Titel des Versuchs}
\newcommand{\datum}{TT.MM.JJJJ}

% Angaben zur Person / Gruppe
\newcommand{\name}{Vorname Nachname}
\newcommand{\gruppe}{XX}
\newcommand{\Kurs}{Vormittags / Nachmittags}

% Hilfsbefehl für einfache Abbildungen:
% \abbildung{Breite als Anteil der Textbreite}{Datei}{Bildunterschrift}{label}
\newcommand{\abbildung}[4]{
    \begin{figure}[ht]
        \centering
        \includegraphics[width=#1\textwidth]{#2}
        \caption{#3}
        \label{#4}
    \end{figure}
}

\begin{document}

% Titelseite und Inhaltsverzeichnis
\FrontMatter
% !TeX root = ../main.tex
% Titelseite: Angaben werden in main.tex geändert.

\begin{titlepage}
    \vspace{20mm}
    \begin{center}
        \scalebox{1.6}{\textbf{\Large \versuch}}\par
        \vspace{20mm}
        \scalebox{1.3}{\textbf{\Large PAP 1}}, \cite{Skript} \cite{Joni}\par
        \vspace{10mm}
        \scalebox{1.2}{\datum}\par
    \end{center}
    \vspace{150mm}
    Teilnehmende Person: \textbf{\name}\par
    Gruppe: \gruppe\par
    Kurs: \Kurs\par
\end{titlepage}

\begingroup
    \let\clearpage\relax
    \pagestyle{plain}
    \newpage
    \tableofcontents
\endgroup

% Hauptteil: Diese Kapitel liegen bewusst in einzelnen Dateien.
\MainMatter
\begingroup
    \let\clearpage\relax

    \chapter{Einleitung}
    % !TeX root = ../main.tex
% In diesem Kapitel stehen Motivation, Messmethode und die nötigen Grundlagen.

\section{Motivation}
% Welche physikalische Fragestellung wird untersucht und warum ist sie relevant?
Hier kommt die Motivation hin
\section{Messverfahren}
% Welches Messprinzip wird verwendet? Welche Größen werden direkt bestimmt?

\section{Physikalische Grundlagen}
% Relevante Modelle, Gleichungen und Annahmen erläutern.


    \chapter{Durchführung}
    % !TeX root = ../main.tex
% Versuchsaufbau und Durchführung knapp und nachvollziehbar beschreiben.

\section{Versuchsaufbau}
% Geräte, Schaltung/Anordnung und Messmittel beschreiben.
% Abbildungen können im Ordner fig/ gespeichert werden.

\section{Durchführung und Messwerterfassung}
% Messablauf, Einstellungen, Wiederholungen und direkt gemessene Größen angeben.
% Große Messwerttabellen können in eine separate Datei im Ordner tab/ ausgelagert werden.


    \newpage
    \chapter{Auswertung}
    % !TeX root = ../main.tex
% Die Aufgaben und Unterabschnitte an das jeweilige Versuchsblatt anpassen.

\section{Aufgabe 1}
% Formelansatz, eingesetzte Messwerte, Ergebnis und Unsicherheit dokumentieren.

\section{Aufgabe 2}

\section{Aufgabe 3}

\section{Aufgabe 4}

% Weitere Aufgaben bei Bedarf ergänzen oder nicht benötigte Abschnitte löschen.


    \chapter{Zusammenfassung und Diskussion}
    % !TeX root = ../main.tex
% Zentrale Ergebnisse knapp zusammenfassen und physikalisch einordnen.

\section{Zusammenfassung der Ergebnisse}
% Wichtigste Endergebnisse mit Einheiten und Unsicherheiten nennen.

\section{Diskussion und Fehlerbetrachtung}
% Ergebnisse mit Literatur-/Sollwerten vergleichen, Abweichungen diskutieren
% und wesentliche zufällige sowie systematische Fehlerquellen bewerten.

\endgroup

% Literaturverzeichnis (Quellen in lit.bib ergänzen und im Text zitieren)
\TheBibliography
\bibliographystyle{natdin}
\bibliography{lit}

\end{document}
